\section{Introduction}\label{sec:introduction}
Network models often combine routing decisions with a choice among services
or operating modes. A route-specific adjustment to a service cost gives a
product of an arc flow and a service weight. Only a small fraction of the
possible route--service pairs may occur. Convexifying each product separately
discards restrictions imposed jointly by flow conservation and the common
simplex. The exact joint hull is available through classical simplex-vertex
disaggregation, but its usual formulation copies every arc for every state.
The question studied here is how the network and the retained products can
reduce that expansion, and when the resulting hull admits simple inequalities
in the original coordinates.

The general convex hull is already understood through disjunctive
formulations and the simultaneous convexification framework of
\citet{DavarniaRichardTawarmalani2017}. Its closest network specialization,
\citet{KhademniaDavarnia2025}, provides a complete aggregation-and-projection
framework and explicit tree and forest inequalities. Our aim is to identify
what the equality-flow structure and the pattern of retained products imply
\emph{within} that known hull: how many state coordinates a constructive
formulation retains, when a complete separator can use a finite structural
library, and when nonunit product coefficients are unavoidable.

We consider a bounded equality-flow polytope
$P=\{x:Ax=b,\ 0\le x\le u\}$ on a directed multigraph and a vector
$y\ge0$ with $\sum_jy_j\le1$. Only products $z_{ej}=x_e y_j$ in a
specified set $\Obs$ are retained. We call these retained product coordinates \emph{observations}
and their joint convex hull $H$. A simplex label is an index $j$;
the unused weight $1-\sum_jy_j$ forms an additional residual state.
The original coordinates are the flows $x$, the weights $y$, and the retained
products $z$.

The key distinction is between network size and unobserved circulation
freedom. Circulations separate across the cyclic blocks of the underlying
undirected graph, and labels absent from a block can be merged there. Within
an observed label, a remaining circulation must be supported entirely on
unobserved arcs. For example, on a cycle, one observed arc determines the
whole state circulation. More generally, if the unobserved arcs form a forest
for each block and observed label, no additional state coordinates are
needed, even when the network itself has many cycles. A complementary
reduction applies to a chain of parallel arc pairs with one bypass: all
serial components share a vector of state flows whose dimension depends on
the number of observed labels, rather than on chain length.

These reductions have two uses. They give exact formulations whose size
reflects the information actually present in the model. They also make it
possible to classify when eliminating state variables preserves simple
inequalities. The latter issue is not settled by total unimodularity of a
single network: sharing aggregate flows among several states and retaining
only selected products changes the projected geometry. Our positive and
negative results locate concrete boundaries for unit coefficients, meaning
coefficients in $\{0,\pm1\}$ on flows and products. Simplex coefficients may
depend on the rational network data.

Our main results are as follows.
\begin{enumerate}[leftmargin=*,itemsep=5pt]
\item \emph{An observation-sensitive exact formulation.}
For a cyclic block $B$ and a label $j$ observed there, let $\rho_{Bj}$
be the cycle rank of its unobserved subgraph.
\Cref{thm:observed-compression} gives a formulation with exactly
$\sum_{B,j}\rho_{Bj}$ additional coordinates and unit flow, product,
and auxiliary coefficients in a specified rational row scaling.
This identifies the residual freedom of an explicit construction; it is
not a minimum over all extended formulations. The same number is the
minimum number of additional \emph{individual products} needed to recover
all state products on the ambient circulation space. The contribution is
the graph-based formulation and its coefficient control, building on
standard elimination by multiplied balance equations.

\item \emph{Complete structural separation and recovery procedures.}
For blocks of maximum cycle rank $r$, finite libraries give exact
separation in $2^{O(r^2)}N$ rational arithmetic operations and compact
recovery in $2^{O(r^3)}N$, where
$N=|V|+|E|+m+|\Obs|$, after observation grouping. Neither procedure
calls an online LP. Cycles, theta blocks, and parallel-path blocks admit
simpler interval, support, and transportation-cut descriptions.
A separate finite-library construction applies to flat chains with a fixed
number of observed labels, despite their unbounded cycle rank.
These procedures provide either a valid separating inequality or a
constructive decomposition in the known exact hull; their contribution is
the specified structural bounds and explicit certificates, rather than
polynomial-time separation itself.

\item \emph{Sharp boundaries and sparse obstructions for coefficients.}
The bounded-rank description has coefficient bound $H_r$ from
\eqref{eq:rank-coefficient-bound}; $H_3=2$ is attained by a unit-capacity
$K_4$ with two explicit labels and five products.
On a flat chain, every observation pattern with at most three observed
labels admits a unit flow/product description, and four labels can force
a ratio of two. With growing label count, Fibonacci ratios occur on simple
series--parallel graphs of maximum degree three with only linearly many
observed products. On unrestricted networks, a coordinate-section transfer
of transportation universality gives arbitrary product-coefficient ratios
with two explicit labels and unit network data. These lower bounds concern
actual retained products and persist under addition of affine-hull
equations; they do not follow merely from finding a nonunit aggregation
multiplier or projecting a facet of a larger product space.
\end{enumerate}

To the best of our knowledge, the residual-cycle formulation with the stated
coefficient control, the complete structural libraries with their stated
bounds, and these coefficient thresholds and sparse constructions have not
been established in the prior work discussed in
\cref{subsec:prior-work}. The underlying disaggregation, product elimination,
transportation projection, and support-function methods are established.
\Cref{tab:prior-comparison} separates those ingredients from the precise
results claimed here; \cref{tab:structural-scope} records the graph classes.

\begin{table}[tb]\centering\small
\caption{Structural scope of the positive results. A unit coefficient is in
$\{0,\pm1\}$ on flow and product coordinates; simplex coefficients may contain
the rational network data. Bounds use the stated row scaling, before any
clearing of simplex denominators.}\label{tab:structural-scope}
\begin{tabular}{>{\raggedright\arraybackslash}p{.25\linewidth}
                >{\raggedright\arraybackslash}p{.44\linewidth}
                >{\raggedright\arraybackslash}p{.20\linewidth}}\toprule
Structure & Exact representation or procedure & Main locator\\\midrule
Arbitrary equality-flow graph & Observation-sensitive EF with
$\sum_{B,j}\rho_{Bj}$ additional coordinates & \Cref{thm:observed-compression}\\[3pt]
Unobserved arcs form forests in each observed block/state & Unit
original-coordinate formulation & \Cref{cor:forest-hull}\\[3pt]
Cycle or theta blocks & Unit interval/support cuts; compact recovery &
\Cref{thm:theta-hull,prop:theta-recovery}\\[3pt]
Parallel-path blocks & Complete subset cuts; transportation-flow separator &
\Cref{thm:parallel-hull}\\[3pt]
Block cycle rank at most $r$ & Coefficients bounded by $H_r$;
finite support and recovery libraries & \Cref{thm:bounded-rank,thm:rank-recovery}\\[3pt]
Flat chain, $a$ observed labels & Fixed-$a$ profile oracle; unit coefficients
for $a\le3$, sharp at four & \Cref{cor:chain-observed-labels}\\\bottomrule
\end{tabular}
\end{table}

Equality conservation is essential to the circulation arguments. Converting
balance inequalities to equalities with slack arcs is possible, but all graph
assumptions must then be checked on the transformed network. The general
network and block results permit arbitrary orientations, parallel arcs,
loops, zero simplex weights, and capacity degeneracy. The flat-chain results
use the specified directed unit-data topology. They do not assert the same
fixed-label threshold on general nested series--parallel networks.

Exact rational implementations and numerical LP formulations make the
distinction between these outputs concrete. The experiments compare with
full disaggregation, global removal of unused labels, native fixed-weight
state-flow LPs, and elementary one-flow or independent-network reductions
where applicable. The strongest globally merged LP is often fastest.
When every label is observed somewhere but each block sees few labels,
local compression still reduces 7,215 state-flow variables to 495 total
variables, or 367 after observed elimination. In that synthetic control,
initial compression reduces the median total time from about 30 to 10
milliseconds; elimination gives the smallest model but takes longer.
The dedicated exact oracle also loses to a simple LP on some long-chain
queries. These results support model-size and certificate benefits without
a universal runtime claim.

\Cref{sec:foundations,sec:blocks} establish the model and known disaggregation
and gluing facts. \Cref{sec:compression,sec:structured-oracles,sec:bounded-rank}
develop observation-sensitive formulations and graph-structured oracles.
\Cref{sec:universality,sec:sp-growth} prove the coefficient obstructions,
and \cref{sec:fixed-chains} gives the complementary fixed-label chain results.
\Cref{sec:computation} reports implementations, verification, experiments,
and a two-supplier service-transportation interpretation.
